%%%%%%%%%%%%%%%%%%%%%%%%%%%%%%%%%%%%%%%%%%%%%%%%%%%%%%%%%%%%%%%%%%%%%%%%
% Copyright 2004 by Till Tantau <tantau@users.sourceforge.net>.        %
%                                                                      %
% Modified (c) 2025 by Paul Nong-Laolam                                %
% for ESPEC North America, INC. GL Controller Training & Distribution  %
% Contact information: <pnong-laolam@espec.com>                        %
%                                                                      %
% In principle, this file can be redistributed and/or modified under   %
% the terms of the GNU Public License, version 2.                      %
%                                                                      %
% However, this file is supposed to be a template to be modified       %
% for your own needs. For this reason, if you use this file as a       %
% template and not specifically distribute it as part of a another     %
% package/program, I grant the extra permission to freely copy and     %
% modify this file as you see fit and even to delete this copyright    %
% notice.                                                              %
%%%%%%%%%%%%%%%%%%%%%%%%%%%%%%%%%%%%%%%%%%%%%%%%%%%%%%%%%%%%%%%%%%%%%%%%

\documentclass{beamer}
% Replace the \documentclass declaration above
% with the following two lines to typeset your 
% lecture notes as a handout:
%\documentclass{article}
%\usepackage{beamerarticle}

% There are many different themes available for Beamer. A comprehensive
% list with examples is given here:
%
% http://deic.uab.es/~iblanes/beamer_gallery/index_by_theme.html
%
% Generally, to try each of them to make a decision, you can uncomment 
% the one in the following list to see which one you like. 
%\usetheme{AnnArbor}
%\usetheme{Antibes}
%\usetheme{Bergen}
%cis275-ch6-w22.pdf\usetheme{Berkeley}
%\usetheme{Berlin}
%\usetheme{Boadilla}
%\usetheme{boxes}
%\usetheme{CambridgeUS}
%\usetheme{Copenhagen}
%\usetheme{Darmstadt}
%\usetheme{default}
%\usetheme{Frankfurt}
%\usetheme{Goettingen}  % selected for usage
%\usetheme{Hannover}
%\usetheme{Ilmenau}
%\usetheme{JuanLesPins}
%\usetheme{Luebeck}
\usetheme{Madrid}
%\usetheme{Malmoe}
%\usetheme{Marburg}
%\usetheme{Montpellier}
%\usetheme{PaloAlto}
%\usetheme{Pittsburgh}
%\usetheme{Rochester}
%\usetheme{Singapore}
%\usetheme{Szeged}
%\usetheme{Warsaw} % See themes.tex for details 
% menu keys
\usepackage[os=win]{menukeys}        % use key shadow menu boxes  
\renewmenumacro{\keys}[+]{shadowedroundedkeys}

% Figures
% explore
\usepackage{wrapfig}    % wrap text aroudn figures 
\usepackage{caption}    % set control caption for minipage env
% end explore
\usepackage{colortbl}
\usepackage{graphicx}
\usepackage{verbatim} 
\usepackage{xcolor}  % See callpkgs.tex for details 

\title{ESPEC GL Controller}

% A subtitle is optional and this may be deleted
\subtitle{Introducing GL Controller \& GL Chamber\\ (Remote Customer Support Team)}

\author{Paul Nong-Laolam}%\inst{1} } %\and S.~Another\inst{2}}
% - Give the names in the same order as the appear in the paper.
% - Use the \inst{?} command only if the authors have different
%   affiliation.

% original marker 
%\institute[Universities of Somewhere and Elsewhere] % (optional, but mostly needed)
%\institute[Company and Affiliation] % (optional, but mostly needed)
\institute[ESPEC NA, Inc.] % (optional, but mostly needed)
{
  %\inst{1}%
  ESPEC North America, Inc. \\ 
  4141 Central Parkway \\
  Hudsonville, Michigan 49426 (USA) 
  %\and
  %\inst{2}%
  %Department of Theoretical Philosophy\\
  %University of Elsewhere
  }
% - Use the \inst command only if there are several affiliations.
% - Keep it simple, no one is interested in your street address.
\date{12/18/2025}
% - Either use conference name or its abbreviation.
% - Not really informative to the audience, more for people (including
%   yourself) who are reading the slides online

\subject{Software} % \subjectComputer Science}
% This is only inserted into the PDF information catalog. Can be left
% out. 

% If you have a file called "university-logo-filename.xxx", where xxx
% is a graphic format that can be processed by latex or pdflatex,
% resp., then you can add a logo as follows:

% \pgfdeclareimage[height=0.5cm]{university-logo}{university-logo-filename}
% \logo{\pgfuseimage{university-logo}}

% Delete this, if you do not want the table of contents to pop up at
% the beginning of each subsection:

\AtBeginSubsection[]
{
  \begin{frame}<beamer>{Outline}
    \tableofcontents[currentsection,currentsubsection]
  \end{frame}
}

% PLACE PAGE NUMBERING
\addtobeamertemplate{navigation symbols}{}{%
    \usebeamerfont{footline}%
    \usebeamercolor[fg]{footline}%
    \hspace{1em}%
    \insertframenumber/\inserttotalframenumber
}
\setbeamercolor{footline}{fg=blue}
\setbeamerfont{footline}{series=\bfseries} 

%%%%%%%%%%%%%%%%%%%%%%%%%%%%%%%%%%%%%%%%%%%%%%%%%%%%%%%%%%%%%%%%%%%
% Let's get started
% BEGIN MAIN DOCUMENT
\begin{document}

\begin{frame}
  \titlepage
\end{frame}

\begin{frame}{ESPEC GL Controller}
  \tableofcontents
  % You might wish to add the option [pausesections]
\end{frame}
%

%%% MAIN MATTER %%%
\input{main-cs2}
%\include{main-cs2}

%%% END MATTER %%%
% All of the following is optional and typically not needed. 
\appendix
%\section<presentation>*{\appendixname}
%\subsection<presentation>*{References}
\section<presentation>*{References}

\begin{frame}[allowframebreaks]
  \frametitle<presentation>{References}
    
  \begin{thebibliography}{10}
    
  \beamertemplatebookbibitems
  % Start with overview books.
%%%
% SEE this site for help on citing the biblio.
% https://tex.stackexchange.com/questions/356146/place-cited-references-into-the-footline-of-slide
% 

  \bibitem{ESPECUWebKit}
    ESPEC.~U-Web Kit.
    \newblock {\em Reference Manual: ESPEC U-Web Kit}, wiki/bitbucket.com, 2022. \\
    {\url{https://bitbucket.org/especnorthamerica/especweb/wiki/EWC_v3-UWeb-Kit}}

  \bibitem{ESPECGLCOMM2025}
    ESPEC.~GL Communications.
    \newblock {\em User Manual: ESPEC GL Controller Communications Manual}, wiki/bitbucket.com, 2025. \\
    {\url{https://bitbucket.org/especnorthamerica/especweb/wiki/ESPEC_GL-Communications-Manual}}

  \bibitem{ESPECGLMAN2025}
    ESPEC.~GL Controller Software.
    \newblock {\em User Manual: ESPEC GL Controller Software, Version 4.0.0}, wiki/bitbucket.com, 2025. \\
    {\url{https://bitbucket.org/especnorthamerica/especweb/wiki/ESPEC_GL-Controller-Operation-Manual}}

  \bibitem{ESPECGLCCLIB2025}
    ESPEC.~GL Python 3 Library.
    \newblock {\em Link to Chamber Connect Library on GitHub}, espec.com/na/support/software, 2025. \\
    {\url{https://espec.com/na/products/controllers/gl_controller}} \\
    Access {\bf{GL support software}} under the above link.  

  \bibitem{ESPECTSERIESUWEBKIT2023}
    ESPEC.~Tseries Web Kit.
    \newblock {\em Link to UWeb Kit}, espec.com/na/support/software, 2023. \\
    {\url{https://bitbucket.org/especnorthamerica/especweb/wiki/EWC_v3.3-Tseries-UKIT}} \\
    Access {\bf{T-series U Web Kit}} under the above link.  

  \bibitem{ESPECTSERIESMAN2023}
    ESPEC.~Tseries Operation Manual.
    \newblock {\em Link to Online Web Controller and T-series Operation Manual}, espec.com/na/support/software, 2025. \\
    {\url{https://bitbucket.org/especnorthamerica/especweb/wiki/EWC_v3.3-Tseries-manual}} \\
    Access {\bf{T-series Operation Manual}} under the above link.  

  %\bibitem{ESPECGLCCLIB2025}
  %  ESPEC.~GL Python 3 Library.
  %  \newblock {\em Link to Chamber Connect Library on GitHub}, espec.com/na/support/software, 2025. \\
  %  {\url{https://espec.com/na/products/controllers/gl_controller#software}}

%  \beamertemplatearticlebibitems
%  \bibitem{UbuntuUnity2017}
%     Unity desktop ended.
%     \newblock Failed to popualrize or market Unity as a solution to all Ubuntu products, Unity project will be ended in April 2017, said Canonical boss Mark Shuttleworth.  Retrieved from {\url{https://www.bbc.com/news/technology-39490848}}. 

  %  S.~Someone.
  %  \newblock On this and that.
  %  \newblock {\em Journal of This and That}, 2(1):50--100,
  %  2000.
  \end{thebibliography}
\end{frame}

\end{document}